%==================================================================
% Ini adalah abstrak dalam bahasa indonesia 
%==================================================================

%% DILARANG EDIT BAGIAN INI
\clearpage
\phantomsection
\addcontentsline{toc}{chapter}{ABSTRAK}
\begin{center}
    \textbf{\large Abstrak}\\[3em]
\end{center}
%\noindent {\penulis}, "{\judulid}", Tugas Akhir/Skripsi DIV Jurusan Teknik
%{\jurusan} Politeknik Negeri Semarang, di bawah bimbingan Nama {\pembimbingsatu}
%dan {\pembimbingdua},{\bulan}, {\tahun}, Jumlah Halaman.
%%----------------------------------------------------------------

\vspace{0.2cm}

Dalam aplikasi kendali gerak \textit{point-to-point} (P-T-P) berkecepatan dan
berpresisi tinggi di ranah industri, induksi getaran struktural sisa
(\textit{residual vibration}) akibat sentakan mekanis (\textit{jerk}) yang
diskontinu menjadi kendala utama, terutama saat menggunakan profil kecepatan
konvensional seperti \textit{trapezoidal}. Untuk mengatasi masalah tersebut,
penelitian ini mengusulkan implementasi algoritma \textit{S-Curve trajectory
planning} berorde tiga guna menekan osilasi mekanis pada sistem penggerak
\textit{Servo Amplifier} Mitsubishi Melservo-J2-Super. Profil \textit{S-Curve}
mampu menghasilkan transisi gerak yang sangat halus, beban komputasi integrasi
numerik orde tiganya sangat tinggi dan berisiko memicu latensi \textit{pulse
jitter} pada unit pemroses standar. Oleh karena itu, penelitian ini merancang
arsitektur sistem kendali \textit{multitasking} deterministik berbasis
\textit{dual-core microcontroller}. Beban kerja dipisahkan secara mutlak:
\textit{Core 0} didedikasikan untuk komputasi matriks trajektor, sedangkan
\textit{Core 1} diisolasi secara eksklusif untuk mengeksekusi \textit{hardware
timer interrupt} guna membangkitkan pulsa kendali secara presisi. Untuk menjamin
integritas pulsa dari interferensi elektromagnetik, sinyal ditransmisikan
menggunakan \textit{differential line driver}. Kinerja sistem dievaluasi secara
eksperimental dan komparatif terhadap profil konvensional menggunakan \textit{AC
servo training kit} ED4270m berbasis aktuator mekanis presisi tinggi
\textit{linear ball-screw}. Selain mengukur deviasi \textit{positional error}
dan durasi waktu tenang (\textit{settling time}). Hasil
penelitian diharapkan memberikan referensi sistem kendali gerak lingkar tertutup
(\textit{closed-loop}) yang andal, berbiaya rendah, serta terbukti secara
empiris mampu memperpanjang usia pakai komponen mekanis.

\vspace{0.5cm}

\noindent Kata kunci: \katakunci